% ============================================================================
% P4 — Supplementary Material
% Pareto-guided Monte Carlo tree search explores multi-objective trade-offs
% missed by scalar antimalarial generation
% Target: Journal of Cheminformatics (BMC)
% ============================================================================
\documentclass[10pt,a4paper]{article}

\usepackage[margin=2.2cm]{geometry}
\usepackage[T1]{fontenc}
\usepackage[utf8]{inputenc}
\usepackage{lmodern}
\usepackage{amsmath,amssymb}
\usepackage{booktabs}
\usepackage{graphicx}
\usepackage{siunitx}
\usepackage{seqsplit}
\usepackage{xcolor}
\usepackage[colorlinks=true,linkcolor=blue,citecolor=blue,urlcolor=blue]{hyperref}
\usepackage[nameinlink,noabbrev]{cleveref}

\sisetup{
  detect-all           = true,
  group-minimum-digits = 4,
  range-phrase         = {--},
  range-units          = single,
}

\title{\textbf{Supplementary Material}\\[0.4em]
\large Pareto-guided Monte Carlo tree search reveals multi-objective
trade-offs in antimalarial generation}

\author{}
\date{}

\begin{document}
\maketitle

% ============================================================================
\section*{S1\hspace{0.6em}Pareto-front data provenance}
% ============================================================================

The canonical \textbf{pre-activity} merged Pareto front (main-text Table~2) comprises four mutually
non-dominated candidates. The public-activity term belongs to the separate v12 scalar
benchmark and was not used to construct this historical front. Every candidate is traceable to the deposited
per-seed Pareto CSV \path{p4_pareto_seed_N.csv} from which it originated,
and all MPO, SA, RRS and PNS values are identical between the per-seed record
and the merged front (\path{results/pareto/merged_pareto_front.csv}). The only
score that differs is SYBA, by design: the original Pareto-MCTS run returned a
constant zero fallback for every molecule (SYBA was inactive during search, as
recorded in the per-seed CSVs), so SYBA was recomputed post-hoc with the
lich/conda classifier and the recomputed values were used to re-derive the
non-dominated front and its hypervolume. This procedure is documented in the
main-text Methods section and is fully reproducible: the script
\path{scripts/p4_recompute_pareto_syba.py} regenerates the merged CSV and the
log \path{merged_pareto_front_recompute.log} records the raw and normalised
SYBA scores (\num{78.67}, \num{58.13}, \num{-19.25}, \num{25.88}) and the
hypervolume (\num{1.2366}). The verification script
\path{scripts/p4_pareto_provenance_check.py} re-runs all of these checks and
exits with status 0 only when every SMILES and score in the merged front
matches a deposited record.

\cref{tab:sm_pareto} lists the four canonical candidates with their provenance.

\begin{table}[htbp]
  \centering
  \caption{Canonical \textbf{pre-activity} merged Pareto front with per-seed provenance. All four
  points are mutually non-dominated on the active objectives (MPO, SYBA, RRS,
  PNS; SA is constant at \num{3.0} and therefore excluded from the
  hypervolume). SYBA was recomputed post-hoc (main-text Methods); MPO, SA,
  RRS and PNS are unchanged from the per-seed records.}
  \label{tab:sm_pareto}
  \small
  \begin{tabular}{p{0.34\textwidth}ccccccl}
    \toprule
    Candidate (SMILES) & MPO & SYBA & SA & RRS & PNS & seed & origin CSV \\
    \midrule
    \seqsplit{CNC(=O)c1ccc(OC(C)C(C)(C)C)c(OC)c1} & \num{0.910} & \num{1.000} & \num{3.0} & \num{0.211} & \num{1.0} & \num{13} & \path{p4_pareto_seed_13.csv} \\
    \seqsplit{COc1c(C)cc(-c2ccc(S(N)(=O)=O)cc2)cc1C} & \num{0.945} & \num{1.000} & \num{3.0} & \num{0.196} & \num{1.0} & \num{18} & \path{p4_pareto_seed_18.csv} \\
    \seqsplit{CC1CCC(n2cccn2)NC1(c1ccccc1)C1OCCO1} & \num{0.946} & \num{0.021} & \num{3.0} & \num{0.141} & \num{0.0} & \num{5} & \path{p4_pareto_seed_5.csv} \\
    \seqsplit{COc1cc(CO)c(O)c(C)c1C} & \num{0.729} & \num{0.994} & \num{3.0} & \num{0.223} & \num{1.0} & \num{6} & \path{p4_pareto_seed_6.csv} \\
    \bottomrule
  \end{tabular}
\end{table}

% ============================================================================
\section*{S2\hspace{0.6em}Molecular diversity of the four generation methods}
% ============================================================================

The main text reports v12 scalar-reward statistics separately, because the primary
contributions are the honest benchmark ranking and the multi-objective Pareto
front. Molecular diversity is analysed here in the Supplementary Material
using real, deposited molecule sets.

\textbf{Molecule sets.} For each of the four generation methods (MCTS+ScafVAE,
random search, greedy search, genetic algorithm) the best-scoring molecule of
each of the \num{20} independent benchmark seeds was recorded, giving a
deposited set of \num{20} molecules per method under
\path{results/benchmark_molecules_opt_v12/}. These are the v12-activity runs whose scalar
rewards are reported in the main-text benchmark table and in
\Cref{tab:sm_benchmark}. The canonical benchmark uses the screened-optimal
MCTS configuration ($c_{\text{PUCT}}=\num{5.0}$, $\nu=\num{0.01}$,
$T=\num{0.8}$, ScafVAE policy wired into PUCT priors) and a value-preserving
vectorisation of the docking Tanimoto scan; all \num{20} per-seed runs
reproduce deterministically from the deposited code and data. An earlier
deposit could not be reproduced from the committed code because the docking
oracle library state differed at run time; the present molecule sets supersede
that deposit. The method ranking (Random $>$ MCTS $>$ GA $>$ Greedy) is
unchanged. This ranking describes the v12 scalar reward only; it is not a
claim that scalar reward and Pareto-front coverage measure the same property.

\textbf{Descriptors.} ECFP4 fingerprints (Morgan radius 2, \num{2048} bits)
were computed with RDKit. Pairwise Tanimoto dissimilarity is
$1 - \text{Tanimoto}$; the reported value is the mean over all unique pairs
(\num{190} pairs per method). Bemis--Murcko scaffolds were extracted with
RDKit's \texttt{MurckoScaffoldSmiles}; the unique-scaffold fraction is the
number of distinct scaffolds divided by the number of valid molecules.
Two-dimensional MDS coordinates (metric MDS on the Tanimoto dissimilarity
matrix) are deposited for visualisation
(\path{results/diversity/p4_diversity_mds.csv}).

\cref{tab:sm_diversity} reports the real diversity statistics, and \Cref{fig:sm_mds} shows the metric-MDS projection of the four molecule sets.

\begin{figure}[htbp]
  \centering
  \includegraphics[width=0.62\textwidth]{Graphics/scaffold_diversity.png}
  \caption{Metric multidimensional-scaling (MDS) projection of the deposited
  best-in-seed molecule sets (\num{20} molecules per method) on ECFP4 Tanimoto
  dissimilarity. Greedy search collapses to a single point (deterministic local
  optimum); the three stochastic methods occupy distinct, broad regions of
  chemical space. Deposited coordinates:
  \texttt{results/diversity/p4\_diversity\_mds.csv}.}
  \label{fig:sm_mds}
\end{figure}

\begin{table}[htbp]
  \centering
  \caption{Real molecular diversity of the four generation methods computed
  from the deposited molecule sets (\num{20} best-in-seed molecules per
  method). Values are computed with RDKit; deposited CSVs:
  \texttt{p4\_diversity\_metrics.csv} (computed from the public-activity-augmented benchmark molecule sets) and \texttt{p4\_benchmark\_molecules\_seed\_N.csv} under \texttt{results/benchmark\_molecules\_opt\_v12/}.}
  \label{tab:sm_diversity}
  \small
  \begin{tabular}{lcccc}
    \toprule
    Method & $n$ & \shortstack{Mean pairwise\\ Tanimoto dissimilarity} & \shortstack{Unique\\ Bemis--Murcko scaffolds} & \shortstack{Unique scaffold\\ fraction} \\
    \midrule
    MCTS+ScafVAE & \num{20} & \num{0.7732} & \num{19} & \num{0.950} \\
    Random & \num{20} & \num{0.8046} & \num{20} & \num{1.000} \\
    Greedy & \num{20} & \num{0.0000} & \num{1} & \num{0.050} \\
    GA & \num{20} & \num{0.7761} & \num{12} & \num{0.600} \\
    \bottomrule
  \end{tabular}
  \smallskip
  \raggedright\footnotesize Mean pairwise dissimilarity is the mean
  $1-\text{Tanimoto}$ over all unique pairs (\num{190} pairs per method);
  scaffolds are unique Bemis--Murcko skeletons (RDKit); the fraction is
  unique scaffolds divided by valid molecules.
\end{table}

\begin{table}[htbp]
  \centering
  \caption{Canonical re-benchmarked scalar rewards with the public-activity
  reward term (\num{20} seeds; mean, standard deviation, min, max). MCTS uses
  the screened optimal configuration (\texttt{c\_puct=5.0},
  \texttt{virtual\_loss=0.01}, \texttt{T=0.8}, ScafVAE policy wired). Deposited
  CSV: \texttt{p4\_benchmark\_merged.csv} under
  \texttt{results/benchmark\_molecules\_opt\_v12/}.}
  \label{tab:sm_benchmark}
  \begin{tabular}{lccccc}
    \toprule
    Method & Mean reward & Std & Min & Max & Mean time (s) \\
    \midrule
    Random & \num{0.6724} & \num{0.0056} & \num{0.6616} & \num{0.6856} & \num{50.4} \\
    MCTS+ScafVAE & \num{0.6649} & \num{0.0068} & \num{0.6554} & \num{0.6779} & \num{92.1} \\
    Greedy & \num{0.4278} & \num{0.0000} & \num{0.4278} & \num{0.4278} & \num{1.9} \\
    GA & \num{0.6453} & \num{0.0124} & \num{0.6178} & \num{0.6715} & \num{2.0} \\
    \bottomrule
  \end{tabular}
\end{table}

Because greedy search is deterministic and returns the same local optimum across all \num{20} seeds
(std identically zero in the canonical benchmark), its \num{20} per-seed
outputs coincide, yielding zero pairwise diversity and a single unique
Bemis--Murcko scaffold by construction (\cref{tab:sm_diversity}). The
three stochastic methods generate structurally distinct outputs per seed:
random search achieves full scaffold coverage (\num{20} unique scaffolds from
\num{20} molecules), while MCTS and GA return \num{19} and \num{12} unique
scaffolds respectively. Mean pairwise Tanimoto dissimilarity is high
(\num{0.7761}--\num{0.8046}) for the three stochastic methods. Scaffold
coverage nevertheless differs across methods (GA \num{12}/\num{20}, MCTS
\num{19}/\num{20}, Random \num{20}/\num{20}), so the diversity result is
reported descriptively rather than as evidence of uniformly broad exploration.

% ============================================================================
\section*{S3\hspace{0.6em}Reproducibility}
% ============================================================================

All deposited CSVs, the benchmark, Pareto and diversity analysis scripts, and
the provenance-check script will be listed in the Zenodo manifest. The canonical
benchmark is produced by \texttt{p4\_mcts\_benchmark.py} with best-SMILES
recording (no change to reward semantics) and a value-preserving vectorisation
of the docking Tanimoto scan (verified identical on held-out molecules). The
\num{20} per-seed rewards in
\path{results/benchmark_molecules_opt_v12/p4_benchmark_merged.csv} reproduce
deterministically from the deposited code and data. Diversity statistics are
computed by \path{scripts/p4_compute_diversity.py} from the deposited molecule
sets. Pareto-front provenance is verified by
\path{scripts/p4_pareto_provenance_check.py}, which exits with status 0 only
when every SMILES and score in the merged front matches a deposited per-seed
record.

\section*{S4\hspace{0.6em}Selection mechanism in Pareto-guided MCTS: relation to ParetoDrug and Mothra}
% ============================================================================

To position the selection scheme used here against the closest Pareto-MCTS
generators, we state the differences explicitly.

\paragraph{ParetoDrug (Yang et al., Commun.\ Biol.\ 7, 1074, 2024;
doi:10.1038/s42003-024-06746-w).} ParetoDrug selects the next
atom symbol with \emph{ParetoPUCT}: it computes a \emph{vector} selection score
per child, retains the children that are non-dominated in that score space, and
draws the next action uniformly at random from that Pareto front of children.
The multi-objective character therefore enters \emph{inside} the tree-selection
criterion.

\paragraph{Mothra (Suzuki et al., J.~Chem.\ Inf.\ Model.\ 64, 7291--7302,
2024; doi:10.1021/acs.jcim.4c00759).} Mothra back-propagates a
\emph{multidimensional reward vector} and ranks Pareto-front nodes during
selection using a hypervolume indicator with a projected-distance penalty
(UCB-style bound over the vectorial reward). Again, the Pareto treatment is
carried through the selection and back-propagation of the search tree.

\paragraph{This work.} The scalar four-method benchmark uses
\texttt{compute\_mpo\_reward} with the v12 fixed weights ($w_{\text{MPO}}=\num{0.36}$,
$w_{\text{docking}}=\num{0.315}$, $w_{\text{SYBA}}=\num{0.135}$,
$w_{\text{SA}}=\num{0.09}$, $w_{\text{activity}}=\num{0.10}$); it is a
public-activity-augmented scalar-reward comparison and is reported
as such. In the Pareto-MCTS search, the objective \emph{vector} is stored for
the front, while the tree selection uses a \emph{scalar proxy}
$Q(s,a)$ obtained by summing the min--max normalised objective components
across the four active objectives (PUCT). That is, in contrast to ParetoDrug
and Mothra, the multi-objective character enters through the \emph{front
maintenance} (non-dominated pool per rollout, merged across seeds), not through
a vectorial or Pareto-ranked selection step. This separation is intentional: it
lets us attribute the returned front to the front-keeping mechanism rather than
to a particular vectorial selection heuristic, and it keeps the scalar benchmark
and the Pareto search on a common PUCT backbone. The consequence is that the
headline result---a Pareto front spanning the potency--accessibility
trade-off---is obtained even though tree selection itself is scalar, which is a
conservative test of the front-keeping contribution.

\paragraph{Empirical check.} We directly tested whether this separation costs
anything by adding a ParetoPUCT-style selection variant that, at every node,
restricts the candidates to the non-dominated children (on the mean
multi-objective vectors) before applying the PUCT score
(\texttt{selection\_mode="pareto"} in \texttt{p4\_mcts\_pareto.py}). At equal
budget (\num{700} iterations, \num{5} seeds) the merged fronts are
indistinguishable under a common normalisation (hypervolume \num{15.56} for the
scalar-proxy rule vs \num{15.49} for the Pareto-only rule, $\Delta \approx
\num{0.5}\%$). This confirms that the scalar proxy is not a bottleneck: the
returned front is attributable to the global non-dominated pool and the
policy-guided exploration, both shared by the two variants.

\end{document}
